\chapter{Veri Guvenligi}

\section*{Giris}
Veri guvenligi, verinin olusturma, saklama, aktarim, kullanim ve imha asamalarinin tamami boyunca gizlilik, butunluk ve erisilebilirlik ilkelerini korumayi hedefler.

\section{Kriptografi ve Sifreleme Temelleri}
\subsection{Blok (Block) ve Akis (Stream) Sifreleme Mimarileri}
Blok sifreleme dosya/veritabani korumasinda, akis sifreleme gercek zamanli iletim senaryolarinda one cikar.

\subsection{Simetrik (AES, 3DES) ve Asimetrik (RSA, Diffie-Hellman) Algoritmalar}
Modern sistemlerde simetrik + asimetrik hibrit model kullanilir: anahtar degisimi asimetrik, veri tasima simetrik.

\subsection{Kriptografik Ozetleme (Hash Fonksiyonlari: MD5, SHA1) ve Carpisma Analizi}
MD5 ve SHA1 butunluk kontrolu icin zayif kabul edilir; yeni tasarimlarda SHA-256/SHA-3 tercih edilmelidir.

\section{Veri Yasam Dongusu ve Siniflandirma}
\subsection{Duragan (At Rest), Hareket Halindeki (In Transit) ve Kullanimdaki (In Use) Veri}
Her veri durumu icin ayri kontrol seti tanimlanmalidir: disk sifreleme, TLS/mTLS, bellek izolasyonu ve erisim denetimleri.

\section{Veri Sizintisi Onleme (DLP)}
\subsection{Ag ve Uc Nokta Tabanli DLP Sistemleri}
DLP politikasi; e-posta, web cikisi, dosya paylasimi ve uc nokta ajanlarinda tutarli uygulanmalidir.

\section{Yedekleme Stratejileri}
\subsection{3-2-1 Kurali ve Degistirilemez (Immutable) Yedekler}
Fidye yazilimi dayanikliligi icin immutable yedek, offline kopya ve geri donus testleri kritik oneme sahiptir.

\Ipucu{Yedek varligi kadar, geri donus testinin periyodik yapilmasi da KPI olarak izlenmelidir.}
